% ---------------------------------------------------------------------
%  CHƯƠNG 5 — KẾT QUẢ VÀ ĐÁNH GIÁ
% ---------------------------------------------------------------------
\chapter{Kết quả và đánh giá}
\label{ch:ketqua}

\begin{muctieu}
Chương này đối chiếu các mục tiêu của đồ án, xác định độ đo đánh giá và nêu
những giới hạn có thể khẳng định từ phạm vi đã chọn.
\end{muctieu}

\section{Kết quả đạt được}

Thư mục \code{AI\_Fashion\_Project} cung cấp notebook, checkpoint, cấu hình
và nhật ký cho các lần chạy V1, V2 và V3. Nhờ đó có thể xác nhận thành phần
nhận diện đã được huấn luyện theo bài toán phát hiện đối tượng 21 lớp và đã
lưu các checkpoint \code{best.pt}. V2 có kết quả trên tập test riêng; V3 có
kết quả validation và các biểu đồ do Ultralytics sinh ra.

Mục tiêu ``độ chính xác 80\%'' ban đầu không tương ứng trực tiếp với chỉ số
của bộ phát hiện. Chỉ số tốt nhất ghi nhận cho V3 full là mAP@0,5:0,95 bằng
36,69\% trên validation, nên chưa có bằng chứng để tuyên bố đạt mục tiêu
80\%. Bộ kiểm thử backend đạt 104/104; bộ Flutter có 139 test đạt và 6 test
widget lỗi. Phạm vi thực nghiệm hiện tại chưa bao gồm phép đo thời gian API
trên thiết bị Android thật.

\section{Độ đo đánh giá}

Đối với bộ phát hiện đối tượng, precision và recall được tính theo
Công thức~\eqref{eq:precision-recall}, trong đó TP, FP và FN lần lượt là dự
đoán đúng, dự đoán sai và đối tượng bị bỏ sót:
\begin{equation}
\mathrm{Precision}=\frac{TP}{TP+FP},\qquad
\mathrm{Recall}=\frac{TP}{TP+FN}.
\label{eq:precision-recall}
\end{equation}

Average Precision là diện tích dưới đường cong precision--recall của một
lớp. mAP là trung bình AP trên các lớp; báo cáo dùng mAP@0,5 và
mAP@0,5:0,95 theo đúng tệp \code{results.csv} của Ultralytics.

Ngoài chất lượng nhận diện, thời gian từ lúc gửi yêu cầu đến khi nhận kết quả
là độ đo hiệu năng phù hợp cho hệ thống. Chức năng gợi ý có thể được đánh giá
qua tỉ lệ phương án hợp lệ, số kết quả trùng và đánh giá thủ công của người
dùng; chức năng thử đồ 2D được xem xét theo khả năng thao tác và độ ổn định
của lớp ảnh. Các độ đo này chưa được thu thập thành một bộ số liệu định lượng
trong phạm vi thực nghiệm của báo cáo.

Bảng~\ref{tab:ketqua} tổng hợp trực tiếp các tệp \code{results.csv} và đầu
ra đánh giá trong notebook. Các dòng không cùng tập dữ liệu, vì vậy chỉ phản
ánh tiến trình thực nghiệm và không được dùng như một so sánh tuyệt đối giữa
các kiến trúc.

\begin{table}[H]
\centering
\caption{Kết quả được ghi nhận ở các phiên bản YOLO.}
\label{tab:ketqua}
\footnotesize
\begin{tabularx}{\textwidth}{>{\bfseries\color{daunavy}}l >{\raggedright\arraybackslash}X c c c c}
\toprule
\rowcolor{gray100}\bfseries\color{daunavy}Lần chạy & \bfseries\color{daumaroon}Tập đánh giá & \bfseries\color{daunavy}Precision & \bfseries\color{daunavy}Recall & \bfseries\color{daumaroon}mAP@0,5 & \bfseries\color{daumaroon}mAP@0,5:0,95 \\
\midrule
V1, epoch 1 & Val. & 75,55\% & 26,56\% & 23,41\% & 15,98\% \\
V2, checkpoint tốt nhất & test V2 & 54,46\% & 39,95\% & 42,33\% & 30,01\% \\
V3 full, epoch 19 & Val. & 60,31\% & 51,03\% & 50,11\% & 36,69\% \\
V3, 5 epoch tiếp tục & Val. & 55,18\% & 49,19\% & 49,53\% & 36,10\% \\
\bottomrule
\end{tabularx}
\end{table}

Đường cong loss và các độ đo của năm epoch tiếp tục được trình bày ở
Hình~\ref{fig:v3-results}. Ma trận nhầm lẫn chuẩn hoá tương ứng ở
Hình~\ref{fig:v3-confusion}, còn một lô dự đoán validation được minh hoạ ở
Hình~\ref{fig:v3-predictions}.

\begin{figure}[H]
\centering
\includegraphics[width=\textwidth]{ai-training/v3-results.png}
\caption{Loss và độ đo của năm epoch fine-tune cuối ở V3.}
\label{fig:v3-results}
\end{figure}

\begin{figure}[H]
\centering
\includegraphics[width=0.92\textwidth]{ai-training/v3-confusion-matrix-normalized.png}
\caption{Ma trận nhầm lẫn chuẩn hoá của V3 trên tập validation.}
\label{fig:v3-confusion}
\end{figure}

\begin{figure}[H]
\centering
\includegraphics[width=\textwidth]{ai-training/v3-validation-predictions.jpg}
\caption{Một lô dự đoán của V3 trên tập validation.}
\label{fig:v3-predictions}
\end{figure}

Kết quả V3 trong báo cáo được lấy từ tập validation vì notebook chưa lưu đầu
ra đánh giá trên tập test kết hợp. Do đó, Bảng~\ref{tab:ketqua} chỉ phản ánh
các kết quả đã được ghi nhận và không suy rộng sang hiệu năng trên tập test
hoặc thời gian phản hồi của thiết bị triển khai.

\section{Bàn luận}

V2 cải thiện rõ mAP so với kết quả V1 được lưu, nhưng V1 mới chỉ có một epoch
trong \code{results.csv} nên phép so sánh này không đại diện cho tiềm năng
của cấu hình 60 epoch. Với V3 full, mAP@0,5:0,95 tốt nhất đạt 36,69\% ở epoch
19. Năm epoch tiếp tục làm train loss giảm, nhưng mAP không vượt kết quả tốt
nhất trước đó; điều này cho thấy việc kéo dài huấn luyện từ \code{last.pt}
không tự động cải thiện khả năng tổng quát hoá.

Ma trận ở Hình~\ref{fig:v3-confusion} cho thấy một số lớp có đường chéo cao,
chẳng hạn hat khoảng 0,86, shoe và sunglass khoảng 0,88. Ngược lại, coat,
belt, trousers, shoesClassic và shoesSports có tỉ lệ đúng thấp; glasses và
zipper có quá ít mẫu validation để đưa ra kết luận ổn định. Hình
~\ref{fig:v3-predictions} cũng cho thấy mô hình có thể phát hiện đúng áo,
giày và quần trên ảnh sản phẩm, nhưng vẫn sinh hộp sai trên ảnh quảng cáo,
ảnh ghép và trường hợp có nhiều món đồ.

Với gợi ý phối đồ, kết quả do Gemini sinh từ hồ sơ và bối cảnh, không phải
thuật toán lấy tối đa ba bộ đồ trong tủ đồ. Báo cáo mới xác nhận luồng xử lý
và kết quả minh hoạ, chưa có khảo sát người dùng để định lượng tính phù hợp
và tính ổn định của phản hồi. Với thử đồ, các test tự động đã kiểm tra ánh xạ
và biến đổi; sai lệch vị trí trên camera thật khi tư thế, ánh sáng và khoảng
cách thay đổi chưa được lượng hoá.

\section{Hạn chế}

Phạm vi ở Chương~\ref{ch:modau} cho thấy hệ thống chỉ nhận diện một trang
phục chính và chưa xác định thương hiệu, chất liệu hoặc kích cỡ chính xác.
Chức năng phân tích dựa trên metadata YOLO và bối cảnh; chức năng Gợi ý không
dùng ảnh, không dùng tủ đồ và phụ thuộc vào Gemini. Hệ thống chưa học sở thích
cá nhân từ phản hồi dài hạn và chưa có bộ dữ liệu chuẩn để chấm chất lượng tư
vấn.

Thử đồ bằng lớp ảnh 2D không tái tạo độ rủ, nếp gấp, che khuất, ánh sáng hay
biến dạng theo cơ thể. Người dùng vẫn phải điều chỉnh vị trí, kích thước và
góc xoay. Ứng dụng chủ yếu hướng tới Android, chưa được mô tả như một hệ
thống thương mại và chưa tích hợp thanh toán.

Dữ liệu 21 lớp mất cân bằng mạnh: trong validation V3, shirt và tshirt có
hơn 1.500 đối tượng mỗi lớp, trong khi boat và zipper chỉ có hai đối tượng,
glasses có chín đối tượng. Một số nhãn của bộ dữ liệu cũ trộn detection với
segmentation hoặc có toạ độ ngoài miền chuẩn hoá; Ultralytics đã bỏ qua 148
ảnh train và 20 ảnh validation. Ngoài ra, các ánh xạ
\code{Rompers\_Jumpsuits} sang dress và \code{Suiting} sang jacket chỉ là
xấp xỉ ngữ nghĩa.

Đánh giá hiện tại chưa có benchmark thời gian phản hồi trên thiết bị thật.
Sáu test widget Flutter chưa đồng bộ với giao diện hiện hành, và quy tắc bảo
mật Firestore mới được đối chiếu ở mức cấu hình, chưa có biên bản kiểm thử
chéo bằng hai tài khoản độc lập. Đây là các giới hạn của kết quả kiểm thử,
không được tính là chức năng đã xác nhận hoàn toàn.
